\section{Étape 4 -- Ajout et configuration d'un équipement sur le réseau}
\label{secEtapeAjoutEquipement}

Une fois le réseau du coupleur configuré (étape~3), chaque équipement à
scruter sur ce réseau doit être ajouté et paramétré individuellement. Cette
étape est illustrée ici pour l'ajout du contrôleur PFC200 sur le réseau
\texttt{NOC\_ETHIP\_CONVOYEUR}.

\subsection{Ajout de l'équipement}

Dans le navigateur de DTM, effectuer un clic droit sur le réseau, puis choisir l'option
\emph{Ajouter}. Sélectionner, dans la liste des équipements proposés (issue
des fichiers EDS importés à l'étape~1), celui correspondant à l'équipement
recherché, puis valider avec le bouton \emph{Ajouter DTM}. Donner à ce DTM
un nom explicite (par exemple \texttt{guichet} pour le PFC200).

\subsection{Paramétrage du service d'échange}

Sélectionner l'équipement ajouté dans le navigateur de DTM et ouvrir sa
configuration (double clic). Adapter le service d'échange (\emph{Exclusive
Owner}) aux besoins de l'application :

\begin{itemize}
    \item \textbf{RPI} (\emph{Requested Packet Interval}) : choisir une
    période compatible avec la périodicité de scrutation requise par le
    procédé (déterminée par ailleurs à partir de la dynamique des
    capteurs) ;
    \item \textbf{Taille des entrées} (rubrique \emph{Entrée T->O}) :
    nombre d'octets produits par l'équipement, tel que défini par son EDS
    ou par le besoin applicatif ;
    \item \textbf{Taille des sorties} (rubrique \emph{Sortie O->T}) :
    nombre d'octets consommés par l'équipement, de la même manière.
\end{itemize}

Valider ces réglages avec le bouton \emph{Appliquer}, puis fermer la
fenêtre.

\subsection{Adressage de l'équipement}

Retourner dans le navigateur de DTM, et ouvrire la configuration du NOC par un double clic. 
Dans la liste des équipements du réseau, sélectionner l'équipement ajouté,
puis :

\begin{itemize}
    \item dans l'\textbf{onglet Paramétrage de l'adresse} : remplacer
    l'adresse IP proposée par défaut par l'adresse IP réelle de
    l'équipement, conforme au plan d'adressage de l'installation ;
    \item dans l'\textbf{onglet Propriétés} : choisir, au niveau de la
    \emph{Gestion des items}, le \textbf{mode d'importation manuel}. Ce
    mode permet de définir des symboles utilisateur pour nommer
    individuellement chacune des données produites et consommées, plutôt
    que de conserver les noms génériques importés depuis l'EDS.
\end{itemize}

\subsection{Nommage des données échangées}

En cliquant sur \emph{Items}, il devient possible de donner un nom à
chacun des bits produits (onglet \emph{Entrée (bit)}) et consommés (onglet
\emph{Sortie (bit)}). Ces noms doivent correspondre aux noms des champs des
types de données utilisateur qui seront ensuite exploités dans le programme
(par exemple les types \texttt{T\_guichet\_IN} et \texttt{T\_guichet\_OUT}
pour le PFC200).

\UPSTIinfo[Convention de préfixage]{
On adopte le préfixe \texttt{m\_nix} pour les entrées (données produites,
booléennes) et le préfixe \texttt{m\_nqx} pour les sorties (données
consommées, booléennes), conformément à la description de ces données
donnée au paragraphe consacré à la cartographie mémoire de l'équipement
(cf. document \texttt{02b\_donnees\_echangees}).
}

Valider ce nommage avec le bouton \emph{Appliquer}.
